\documentclass[11pt]{article}
\usepackage[margin=1in]{geometry}
\usepackage{amsmath,amssymb,amsthm,booktabs,microtype}
\usepackage[hidelinks]{hyperref}

\newcommand{\eps}{\{\!-1,1\!\}}
\newcommand{\F}{F}
\newcommand{\M}{M}
\newcommand{\PP}{\mathbb P}
\newcommand{\R}{\mathbb R}
\newcommand{\T}{\mathsf T}
\DeclareMathOperator{\dist}{dist}

\newtheorem{theorem}{Theorem}[section]
\newtheorem{proposition}[theorem]{Proposition}
\newtheorem{lemma}[theorem]{Lemma}
\newtheorem{corollary}[theorem]{Corollary}
\theoremstyle{remark}
\newtheorem{remark}[theorem]{Remark}

\title{Finite structure in a min--max quadratic sign problem}
\author{Rob Sneiderman}
\date{August 2026}

\begin{document}
\maketitle

\begin{abstract}
For a symmetric zero-diagonal sign matrix, let $M(A)$ be the largest absolute
value of its quadratic form on the Boolean cube, and put
$F(n)=\min_A M(A)$.  We give an auditable account of the exact values through
order fifteen, including explicit computational trust boundaries.  We also
prove an exact energy-weighted covering formula for one-vertex extension,
record finite non-heredity of optimizers, and give a sharp completion theorem
for one fixed oriented-Hadamard frame which implies $F(16)\le30$.  A
completion search over the banked order-14 catalogue then reduces the
remaining decision $F(16)\in\{28,30\}$ to a rigidity question about
order-15 minimizers, alongside a complete census of the non-rigid ones.
The asymptotic question whether $F(n)/n^{3/2}$ converges remains open.
\end{abstract}

\begin{center}
\fbox{\parbox{0.91\linewidth}{\textbf{Status.}
This is a focused manuscript draft.  Every unconditional analytic statement
is proved below.  Claims labelled computer-assisted depend on the stated
enumeration boundary.  No result here proves convergence or nonconvergence.}}
\end{center}

\section{Problem and elementary structure}

Let $A=(a_{ij})$ range over symmetric zero-diagonal matrices with
$a_{ij}\in\eps$.  Define
\[
 Q_A(x)=\sum_{1\le i<j\le n}a_{ij}x_ix_j,
 \qquad M(A)=\max_{x\in\eps^n}|Q_A(x)|,
 \qquad F(n)=\min_A M(A).
\]
Switching by $z\in\eps^n$ sends $A$ to $D_zAD_z$ and leaves $M(A)$
unchanged.  Vertex permutation and global negation also preserve $M(A)$.

\begin{lemma}[Puncturing and parity]
For every $n\ge2$,
\[
 F(n)\le F(n+1)\le F(n)+n,
 \qquad
 F(n)\equiv\binom n2\pmod2,
\]
and consequently
\[
 F(n)\ge F(n-1)+((n-1)\bmod2).
\]
\end{lemma}

\begin{proof}
Restrict an order-$(n+1)$ signing to its first $n$ vertices.  For every old
spin $x$, its restricted energy is the average of the two full energies
obtained by setting the new spin to $1$ and $-1$.  This gives
$F(n)\le F(n+1)$.  Extending an optimal order-$n$ signing by arbitrary
incident signs changes every energy by at most $n$, proving the upper bound.
Every energy is a sum of $\binom n2$ signs and therefore has that parity; the
maximum and its minimum over signings have the same parity.  When $n-1$ is
odd, the parities at consecutive orders differ, so monotonicity improves by
one.
\end{proof}

There is also a useful exact coding formulation.  Put $E_n=\binom n2$ and
\[
 D_n=\{(x_ix_j)_{i<j}:x\in\eps^n\}
      \cup\{-(x_ix_j)_{i<j}:x\in\eps^n\}.
\]
If $\rho(D_n)$ denotes its covering radius in the Hamming cube, then
\begin{equation}
 F(n)=E_n-2\rho(D_n).                                      \label{eq:cover}
\end{equation}
Indeed, correlation with a sign word is $E_n-2$ times Hamming distance, and
adjoining the negatives converts absolute correlation to nearest-codeword
distance.  This identity is exact but does not by itself control the
second-order deficit of size $n^{3/2}$.

\section{Exact values through order fifteen}

\begin{theorem}[Exact finite sequence]
The exact values from order two through order fifteen are
\[
\begin{array}{c|rrrrrrrrrrrrrr}
n&2&3&4&5&6&7&8&9&10&11&12&13&14&15\\ \hline
F(n)&1&3&4&4&5&9&10&12&13&17&18&20&21&27.
\end{array}
\]
The lower bounds at orders eleven, thirteen, and fifteen are
computer-assisted with the nauty completeness boundary described below.
\end{theorem}

\begin{proof}[Certificate architecture]
For orders at most ten, switching makes every edge incident with a chosen
root positive.  The remaining signing is an ordinary graph on $n-1$
vertices.  Isomorph-free generation, exact Gray-code energy evaluation, and
independent decoding reproduce
\[
 1,3,4,4,5,9,10,12,13.
\]
All energy and threshold decisions use integers.

For order eleven, the certificate scans all $12{,}005{,}168$ unlabeled
graphs on ten residual vertices.  Exact degree and energy filters leave
$2{,}153{,}606$ eligible classes; none admits maximum at most fifteen.  An
explicit order-eleven signing has maximum seventeen.  Parity excludes
sixteen, hence $F(11)=17$.  The independently checked one-vertex extension
then gives an order-twelve signing with maximum eighteen.  Puncturing and
parity give $F(12)\ge18$.

For order thirteen, eight fail-closed nauty shards together contain exactly
$1{,}018{,}997{,}864$ unlabeled graphs on eleven residual vertices.  The
certificate pins each producer's exit status, byte count, record count, and
SHA-256 stream digest.  Exactly two rooted records survive the order-twelve
threshold; direct evaluation gives internal maximum eighteen and minimum
one-vertex extension value twenty-four for both.  Therefore an order-thirteen
signing of maximum eighteen cannot exist.  Parity excludes nineteen.  Every
thirteen-vertex principal submatrix of the order-fourteen Paley conference
matrix has maximum twenty, so $F(13)=20$.  Monotonicity and parity give
$F(14)\ge21$, while direct evaluation of that conference matrix gives
maximum twenty-one.

For order fifteen, a threshold tower decides the value.  A complete nauty
pass over the same $1{,}018{,}997{,}864$ residual graphs keeps the
$82{,}502{,}142$ order-twelve signings with maximum at most twenty-four.
Two one-vertex expansion levels with canonical deduplication keep the
$282{,}202{,}131$ order-thirteen signings with maximum at most twenty-four
and then the $1{,}313{,}164$ order-fourteen signings with maximum at most
twenty-five.  Every order-fifteen signing with maximum at most twenty-five
restricts to a member of that final level, and a complete scan of all
$2^{13}$ projective extensions of every member finds none with maximum at
most twenty-five.  Parity excludes twenty-six, so $F(15)\ge27$.  Extending
the order-fourteen conference matrix by its best sign row gives maximum
exactly twenty-seven, hence $F(15)=27$.  The final tower level is banked as
a digest-pinned catalogue, and its complete extension scan is replayed
exactly by the public certificate; completeness of the two lower levels is
the recorded tower boundary.

The exhaustive lower certificates trust nauty's completeness and canonical
isomorph-free generation.  They do not infer infeasibility from a solver
timeout.  The explicit upper witnesses and all their complete energy scans
are independently reproducible without that trust boundary.
\end{proof}

\begin{remark}
The order-thirteen scan finds two \emph{rooted records}, not a theorem that
there are only two unrooted optimal switching classes.  Likewise, the
certificate proves that the order-fourteen conference matrix is optimal but
does not by itself prove uniqueness.
\end{remark}

\section{The exact one-vertex state}

For an order-$n$ signing $B$, let $E(B)$ be the smallest maximum obtainable
by adjoining one vertex with arbitrary incident sign vector.  Write
\[
 \PP_n=\eps^n/\{x\sim-x\},
 \qquad d_\pm([u],[v])=\min\{d_H(u,v),d_H(u,-v)\}.
\]
Put $M=M(B)$ and define
\[
 w_B([x])=\frac{M-|Q_B(x)|}{2},
 \qquad
 \rho_{\rm w}(B)=\max_{[b]\in\PP_n}\min_{[x]\in\PP_n}
 \bigl(d_\pm([b],[x])+w_B([x])\bigr).
\]
The weights are nonnegative integers because all order-$n$ energies have the
same parity.

\begin{theorem}[Weighted Bellman identity]
For every order-$n$ signing $B$,
\begin{equation}
 E(B)=M(B)+n-2\rho_{\rm w}(B).                            \label{eq:bellman}
\end{equation}
Equivalently, with
$\delta_{\rm w}(B)=\lfloor n/2\rfloor-\rho_{\rm w}(B)$,
\begin{equation}
 F(n+1)=\min_B\bigl(M(B)+(n\bmod2)+2\delta_{\rm w}(B)\bigr).
                                                               \label{eq:bellman2}
\end{equation}
Only configurations satisfying
$|Q_B(x)|\ge M(B)-2\lfloor n/2\rfloor$ can affect the inner covering
minimum.
\end{theorem}

\begin{proof}
Let $b\in\eps^n$ be the incident signs and $t$ the new spin.  Pairing the two
values of $t$ gives
\[
 \max_{t\in\eps}|Q_B(x)+t\,b\cdot x|
   =|Q_B(x)|+|b\cdot x|.
\]
Since $|b\cdot x|=n-2d_\pm([b],[x])$, the maximum over old spins for fixed
$b$ equals
\[
 M+n-2\min_{[x]\in\PP_n}
       \bigl(w_B([x])+d_\pm([b],[x])\bigr).
\]
Minimizing over $[b]$ proves \eqref{eq:bellman}.  Minimizing over $B$ and
using $n-2\lfloor n/2\rfloor=n\bmod2$ proves \eqref{eq:bellman2}.
If $w_B([x])>\lfloor n/2\rfloor$, that state loses to an exact maximizer at
projective distance at most $\lfloor n/2\rfloor$, proving the final claim.
\end{proof}

The formula explains why exact optimizers need not be valid predecessors.
It also shows why a scalar energy is not a closed state: the spatial
placement of every energy layer in the projective cube matters.

\section{Finite optimizer non-heredity}

For fixed internal blocks $A$ and $B$, define
\[
 J(A,B)=\min_{C\in\eps^{s\times k}}
 M\begin{pmatrix}A&C\\C^{\T}&B\end{pmatrix}.
\]

\begin{proposition}[A complete $2+8$ obstruction]
Among optimal blocks of orders two and eight,
\[
 \min_{M(A)=F(2),\,M(B)=F(8)}J(A,B)=15>F(10)=13.
\]
Consequently no optimal order-ten signing contains an optimal order-eight
principal submatrix.  The smallest internal maximum among order-eight blocks
admitting a $2+8$ completion of maximum thirteen is twelve.
\end{proposition}

This is a finite exhaustive proposition.  The verifier enumerates every
relevant switching-permutation class, every rectangular cross block modulo
the valid symmetries, and every projective full spin.  It establishes an
``internal sacrifice'' phenomenon: minimizing the blocks separately can make
their joint optimum worse.

There is a second exact failure of scalar state.  Two optimal order-nine
classes have the same complete projective absolute-energy histogram
\[
 \#\{|Q_B|=0,4,8,12\}=(60,111,60,25),
\]
so all their scalar partition functions agree at every temperature, but their
minimum one-vertex extension values are thirteen and fifteen.  Their weighted
covering geometry, not their energy multiset, distinguishes them.

\section{A sharp oriented-Hadamard completion}

Let
\[
H=\begin{pmatrix}
1&1&1&1\\1&1&-1&-1\\1&-1&1&-1\\-1&1&1&-1
\end{pmatrix},
\quad
(a_{01},a_{02},a_{03},a_{12},a_{13},a_{23})
=(1,1,1,-1,1,-1).
\]
For arbitrary symmetric zero-diagonal order-four signings
$D_0,D_1,D_2,D_3$, let $L(D_0,D_1,D_2,D_3)$ have diagonal blocks $D_i$,
upper cross block $ij$ equal to $a_{ij}H$, and transpose blocks below.

\begin{theorem}[Sharp completion of the pinned frame]
\begin{equation}
 \min_{D_0,D_1,D_2,D_3}M(L(D_0,D_1,D_2,D_3))=30.          \label{eq:frame30}
\end{equation}
In particular, $F(16)\le30$.
\end{theorem}

\begin{proof}
Put
\[
u_0=(1,1,1,1),\quad u_1=(1,1,1,-1),\quad
u_2=(1,1,-1,1),\quad u_3=(1,1,-1,-1).
\]
Direct multiplication gives the following cross-only energies:
\[
\begin{array}{c|c|r}
& (x_0,x_1,x_2,x_3)&Q_{\rm cross}\\ \hline
1+&(u_0,u_0,-u_0,u_1)&28\\
1-&(u_0,-u_3,-u_0,-u_1)&-28\\
2+&(u_3,u_3,-u_3,u_2)&28\\
2-&(u_3,-u_2,-u_3,-u_2)&-28\\
3+&(u_3,u_2,-u_3,u_2)&28\\
3-&(u_3,-u_0,-u_3,-u_2)&-28.
\end{array}
\]
Within each pair the projective spins on clouds zero, two, and three agree,
so their internal contributions cancel in the comparison.  If the full
maximum were at most twenty-eight, the three comparisons would force
\[
 q_{D_1}(u_0)\le q_{D_1}(u_3)\le q_{D_1}(u_2)
 \le q_{D_1}(u_0).
\]
All three values would be equal.  But
\[
 q_{D_1}(u_0)-q_{D_1}(u_2)
 =2(d_{13}+d_{23}+d_{34})\ne0,
\]
because a sum of three signs cannot vanish.  Order-sixteen energies are even,
so the restricted minimum is at least thirty.

For attainment, list the six internal edges in order
$01,02,03,12,13,23$ and take
\[
 P=(1,1,1,-1,1,-1),\qquad
 R=(-1,-1,1,1,1,-1),\qquad
 (D_0,D_1,D_2,D_3)=(P,R,P,R).
\]
An exhaustive exact scan of all $2^{15}$ projective spins gives maximum
thirty.  Independent direct and blockwise implementations reproduce the
complete energy histogram.
\end{proof}

Both $P$ and $R$ have maximum four and are switching-permutation equivalent.
Nevertheless their ordered response vectors to this frame are
\[
 (2,0,4,-2,0,2,-2,-4),\qquad
 (0,-2,2,-4,-2,0,4,2).
\]
Using one literal common internal block in all four clouds has minimum
thirty-eight.  Thus neither $M(D)$ nor the unframed switching class is a
closed substitution state; orientation relative to the cross frame is
essential.  The theorem is restricted to one frame.  It neither proves
$F(16)=30$ nor supplies an iterable amplification law.

\section{A reduction at order sixteen}

Monotonicity from $F(15)=27$, even order-16 parity, and the completion
theorem of the previous section leave exactly two candidates,
\[
 F(16)\in\{28,30\}.
\]
Neither value is claimed here, but the first case is now tightly
constrained.

\begin{theorem}[Pair-deletion rigidity]
Every order-16 signing $W$ with $M(W)\le28$ has all $\binom{16}{2}=120$
order-14 pair-deletions at maximum exactly $27$.  Consequently every
order-15 deletion of $W$ is a minimizer all of whose own deletions have
maximum $27$.
\end{theorem}

\begin{proof}[Certificate architecture]
Each order-15 deletion of $W$ has odd maximum at most $28$, hence exactly
$27$.  Suppose some pair-deletion $B$ had maximum at most $25$.  Then $B$
lies, up to switching and relabeling, in the complete order-14 catalogue of
the order-15 certificate, and the two rows of $W$ incident to the deleted
pair must each extend $B$ to order-15 maximum exactly $27$.  An exhaustive
scan finds that exactly $163$ of the $1{,}313{,}164$ catalogue members
admit such extension rows, with $316$ rows in all, and a complete check of
every completion (both rows and the mutual coupling sign, over all $2^{15}$
spin states, in exact integer arithmetic) finds none of maximum at most
$28$.  The completion minimum is $30$, attained; the special case in which
the pair-deletion is the order-14 conference matrix bottoms at $32$.
\end{proof}

\begin{theorem}[Census of non-rigid minimizers]
Exactly $30$ rooted canonical classes of order-15 signings have maximum
$27$ and some vertex-deletion of maximum at most $25$; the $316$
extensions above exhaust them.  Each member is re-verified at maximum
exactly $27$.
\end{theorem}

Two remarks bound the claim.  The census counts rooted classes, with the
same caveat as the tower levels: distinct rooted records may coincide
under another choice of root.  And the reduction inherits the catalogue
completeness chain of the order-15 certificate.  What remains open at
order sixteen is precisely the rigid case: an order-15 minimizer all of
whose deletions have maximum $27$ would be needed to build a signing of
maximum $28$, and no such minimizer is currently known.  The completion
minimum $30$ gives an independent second family attaining the
framed-Hadamard ceiling.

\section{Reproducibility and scope}

The repository contains source-built verifiers for all statements above.
The principal commands are:
\begin{verbatim}
python3 verification/research_exact_small_n.py --max-n 10
python3 verification/research_order11_certify.py
python3 verification/research_order13_certify.py --full-stream --jobs 8
python3 verification/research_order15_certify.py --full --jobs 4
python3 verification/research_order16_reduction.py --jobs 4
python3 verification/research_cross_block_composition.py
python3 verification/verify_nonlinear_bellman.py
python3 verification/verify_framed_hadamard_lift_30.py
\end{verbatim}
The order-eleven and order-thirteen complete lower scans require nauty and
are much more expensive than witness verification.  The order-fifteen
certificate replays the deciding tower level from a committed digest-pinned
catalogue; re-deriving the lower tower levels from the committed driver
takes about a day of machine time.  Exact commands, pinned counts, expected
outputs, corruption controls, and the producer trust boundary are recorded
in \texttt{verification/README.md}.

These finite results identify optimizer-sensitive composition as the main
structural difficulty.  They do not determine an asymptotic constant and do
not constitute an answer to the MathOverflow limit question.  The companion
manuscript \texttt{composition\_framework.tex} isolates the remaining
cross-order theorem.

\begin{thebibliography}{9}
\bibitem{mo}
P. Ivanisvili,
\emph{Min-max of a quadratic form of plus-minus ones},
MathOverflow question 413935, 2022.

\bibitem{nauty}
B. D. McKay and A. Piperno,
\emph{Practical graph isomorphism, II},
Journal of Symbolic Computation 60 (2014), 94--112.
\end{thebibliography}

\end{document}
